% !TEX program = xelatex
\documentclass[UTF8,a4paper,zihao=-4]{ctexart}

\usepackage[a4paper,top=2.5cm,bottom=2.5cm,left=2.7cm,right=2.7cm]{geometry}
\usepackage{amsmath,amssymb}
\usepackage{booktabs,tabularx,longtable,array}
\usepackage{graphicx}
\usepackage{xcolor}
\usepackage{enumitem}
\usepackage{caption}
\usepackage{subcaption}
\usepackage{tikz}
\usetikzlibrary{arrows.meta,positioning,calc,fit,shapes.geometric}
\usepackage[hidelinks]{hyperref}
\usepackage[backend=biber,style=numeric,sorting=none,giveninits=true]{biblatex}
\addbibresource{LPR_academic_survey_refs.bib}

\setlength{\parindent}{2em}
\setlength{\parskip}{0.25em}
\linespread{1.18}
\setlist[itemize]{leftmargin=2em,itemsep=0.15em,topsep=0.25em}
\setlist[enumerate]{leftmargin=2.2em,itemsep=0.15em,topsep=0.25em}
\captionsetup{font=small,labelfont=bf,labelsep=quad}

\title{\heiti 学习路径推荐研究综述：从传统优化到强化学习与大模型增强}
\author{唐耿}
\date{2026年7月}

\begin{document}
\maketitle

\begin{abstract}
学习路径推荐（Learning Path Recommendation, LPR）旨在根据学习者历史、当前知识状态和学习目标，生成具有教学顺序约束的学习项目序列。与一般内容推荐不同，LPR 的目标不是最大化点击率，而是提升学习收益、路径合理性和长期知识保持。本文结合已有本地文献整理、学习路径推荐相关论文和补充网络检索结果，对该领域的技术演进进行综述。本文将 CSEAL 视为一个重要分界点：CSEAL 之前的研究多依赖本体、知识图、启发式搜索、协同过滤、序列模型或传统组合优化；CSEAL 之后，研究逐步转向以强化学习、层级强化学习、图神经网络、难度感知和离线训练为核心的序列决策框架。由于 LPR 的状态估计高度依赖知识追踪（Knowledge Tracing, KT），本文进一步讨论了 DKT、GKT、RKT、CKT、DIMKT 等 KT 模型如何为路径推荐提供学生状态、遗忘、难度和过程一致性信息。最后，本文归纳了当前研究在评估体系、在线验证、可解释性、知识结构自动构建、长期学习效果、大模型对齐以及逐题闭环式细粒度推荐方面的开放问题，并提出若干可能的研究方向。
\end{abstract}

\noindent\textbf{关键词：} 学习路径推荐；知识追踪；强化学习；个性化学习；知识图谱；大语言模型

\section{引言}

在线教育平台使学习者能够接触到数量庞大的知识点、视频、习题和课程资源，但资源丰富并不自动等于学习高效。学习者常面临两个直接困难：其一，不清楚当前最需要补足的知识点；其二，即使知道目标知识点，也难以确定合适的学习顺序和练习难度。学习路径推荐正是为解决这一类序列化学习决策问题而提出，其核心任务是：在给定学习者历史交互、目标知识集合和有限学习预算的条件下，推荐一个或多个有序学习项目，使学习者从当前状态更有效地到达目标状态。

从推荐系统角度看，LPR 至少具有四个区别于传统推荐的性质。第一，学习项目之间存在先修依赖，路径顺序本身具有教育意义。第二，学习者状态会随交互动态变化，推荐动作不仅满足当前需求，也会改变未来状态。第三，目标函数应反映学习收益和知识保持，而不是简单的点击、停留或购买行为。第四，推荐需要在有限预算内进行长程规划，具有典型的序列决策特征。因此，LPR 很自然地连接了推荐系统、知识追踪、认知诊断、强化学习和教育学理论。

现有研究大体经历了三个阶段。早期工作以本体、知识图、规则、约束优化和序列推荐为主，强调可解释路径和先修关系，但通常缺乏对学习收益的直接优化。2019 年的 CSEAL 将认知结构、学习者状态追踪和 Actor-Critic 策略学习结合起来，为后续强化学习式 LPR 奠定了基本范式\cite{Liu2019CSEAL}。此后，GEHRL、SRC、DLPR、GEPKSD、LIGHT、KnowLP、UNO 等工作围绕层级决策、集合到序列生成、难度感知、知识状态蒸馏、知识拓扑优化、GraphRAG 建图和离线训练继续推进\cite{Li2023GEHRL,Chen2023SRC,Zhang2024DLPR,Li2024GEPKSD,Yu2025LIGHT,Cheng2025KnowLP,Peng2026UNO}。近两年，GenAL、ReAL、IB-GRPO 和 UniER 等工作又把大语言模型、检索增强代理、指标对齐和统一基准带入该领域\cite{Lv2025GenAL,Lv2025ReAL,Wang2026IBGRPO,Cheng2026UniER}。

本文的组织如下：第二节介绍 LPR 在教育辅导系统、在线练习平台和 MOOC 场景中的应用形态；第三节总结代表性 LPR 论文引言中的问题谱系；第四节给出 LPR 的形式化定义与评价目标；第五节回顾 CSEAL 之前的传统方法；第六节讨论 KT、KSS/KES 与学生模拟器在 RL-LPR 中的作用；第七节以 CSEAL 之后的强化学习方法为重点梳理主流路线；第八节总结大模型与知识图谱增强趋势；第九节介绍常用数据集与基准选择逻辑；第十节比较实验评估和指标；第十一节讨论开放问题与未来方向。

\section{应用场景与教育辅导系统}

现有 LPR 论文的引言通常并不只是把学习路径推荐描述成一个抽象算法问题，而是把它放在在线教育平台、智能辅导系统、MOOC 平台和练习推荐系统中讨论。原因在于，不同应用系统改变了推荐动作、反馈信号和优化目标的形态：在题库型辅导系统中，动作往往是“下一道题”或“下一组练习”；在课程平台中，动作可能是“下一门课”“下一个章节”或“一个完整学习单元”；在智能教师代理中，动作还可能包含解释、提示和学习策略建议。因此，应用背景不是论文引言中的修饰性内容，而是决定 LPR 问题建模粒度的重要前提。

从数学上看，一个 LPR 应用场景可以粗略表示为
\begin{equation}
\Omega=(\mathcal{A},\mathcal{O},\rho,H),
\end{equation}
其中 $\mathcal{A}$ 表示可推荐的动作空间，$\mathcal{O}$ 表示系统可观测反馈，$\rho$ 表示评价学习收益的奖励或指标，$H$ 表示规划时间范围。若 $\mathcal{A}$ 是习题集合，$\mathcal{O}$ 通常包括对错、用时和尝试次数，系统可以进行较密集的状态更新；若 $\mathcal{A}$ 是课程或章节集合，$\mathcal{O}$ 往往变成完成率、测验成绩和退课行为，反馈更稀疏，规划周期更长。由此可见，同样称为“学习路径推荐”，其算法困难会随应用系统发生显著变化。

\subsection{在线题库与智能辅导系统}

CSEAL 的引言把 adaptive learning 放在 Khan Academy、Junyi Academy 等在线教育系统背景下讨论，这类系统的核心特征是学习者持续完成练习，系统根据答题表现估计其认知状态，再给出后续学习项目\cite{Liu2019CSEAL}。在这种环境中，LPR 与传统“内容推荐”的差异非常明显：推荐一道题不只是预测学习者是否喜欢这道题，而是希望这道题能暴露当前知识缺口、巩固薄弱知识点，或为后续更高层知识做准备。

这类教育辅导系统通常形成一个闭环：
\begin{equation}
\mathcal{H}_{u,t}\rightarrow \hat{\mathbf{s}}_{u,t}\rightarrow a_t\rightarrow r_t\rightarrow \mathcal{H}_{u,t+1},
\end{equation}
其中 $\mathcal{H}_{u,t}$ 是学习者历史，$\hat{\mathbf{s}}_{u,t}$ 是系统估计的知识状态，$a_t$ 是被推荐的学习项目，$r_t$ 是学习反馈。这个闭环解释了为什么 CSEAL 需要同时考虑学习者知识水平和学习项目之间的知识结构：如果只看答题正确率，系统可能推荐当前容易得分但教育价值较低的题；如果只看知识图谱，系统又可能忽略学习者的个体差异。CSEAL 中用于模拟学习者状态转移的 KSS 基于 IRT 思想而不是 KT，KES 则进一步评估推荐路径带来的学习效果；这说明早期 RL-LPR 已经意识到真实在线系统中很难频繁进行大规模线上试错，因此必须构造可用于离线训练和评估的学生环境\cite{Liu2019CSEAL}。

\subsection{MOOC 与课程级学习平台}

另一类常见应用来自 MOOC 和课程平台。LIGHT 的引言把 LPR 放在 Coursera、XuetangX 等在线教育平台快速发展的大背景下，强调平台中学习资源数量巨大，学习者需要系统帮助其形成有序、可执行的学习路线\cite{Yu2025LIGHT}。与题库型系统相比，MOOC 场景的动作粒度更粗：系统推荐的可能是一门课程、一个章节、一组视频或一个知识模块。学习者完成这些动作需要更长时间，反馈也更容易受到兴趣、时间安排和外部动机影响。

这种场景使 LPR 更接近长期规划问题。面向 MOOC 的课程推荐工作采用层级强化学习，正是因为课程选择既包含高层目标规划，也包含低层资源选择\cite{Zhang2019MOOC}。如果直接在所有课程上进行一步式推荐，动作空间会很大，且难以刻画课程之间的先修关系和目标依赖；层级结构把“先学哪个方向”与“在该方向中选择什么资源”拆开，能够降低搜索难度，也更符合课程学习的组织方式。SRC 所关注的完整路径生成同样对应这类需求：在某些课程组织或阶段性学习计划中，系统需要一次性给出较完整的路径，而不是依赖每一步即时反馈逐渐修正\cite{Chen2023SRC}。

\subsection{练习推荐、考试准备与短期学习计划}

在练习推荐系统中，LPR 往往服务于更具体的目标，例如单元复习、作业辅导、考试准备或薄弱知识点强化。此时系统的任务不是单纯安排知识点顺序，还要决定每个知识点上使用哪些习题、使用多少题、难度如何变化。DLPR 的引言之所以强调 item difficulty，是因为真实练习路径中“知识点顺序正确”并不等于“练习体验合理”：若连续给出远超能力水平的题，学习者可能产生挫败；若长期停留在过易题目上，又会浪费学习预算\cite{Zhang2024DLPR}。

UNO 和 UniER 进一步说明，习题级 LPR 的应用需求正在从单步 exercise recommendation 扩展到 path-level exercise recommendation\cite{Peng2026UNO,Cheng2026UniER}。在单步推荐中，系统只需判断下一题是否合适；在路径级推荐中，系统还要考虑整条练习序列是否覆盖目标知识、是否难度平滑、是否避免重复冗余，以及是否能在有限题量内产生可测量的学习收益。这个差异很重要，因为许多教育辅导系统真实面对的是“给学生安排一段练习计划”，而不是孤立地推荐一道题。

\subsection{生成式教师代理与解释型辅导}

近期 LPR 论文的引言开始把大语言模型和教师代理纳入应用背景。GenAL 和 ReAL 分别把 LLM 用作生成式自适应学习代理和检索增强教师模拟器，试图让系统不只推荐学习项目，还能利用题目文本、知识点语义和教学经验生成更接近教师的引导\cite{Lv2025GenAL,Lv2025ReAL}。IB-GRPO 则进一步指出，仅让 LLM 生成看似合理的路径并不充分，推荐结果还需要与教育目标、学习收益指标和路径约束对齐\cite{Wang2026IBGRPO}。

这一应用形态改变了 LPR 的输出边界。传统 LPR 输出通常是项目序列 $\mathcal{P}=(a_1,\ldots,a_T)$，而教师代理式系统可能输出
\begin{equation}
\mathcal{Y}=\{(a_t,e_t,h_t)\}_{t=1}^{T},
\end{equation}
其中 $e_t$ 是解释，$h_t$ 是提示或辅导话术。这样做的潜在价值在于，LLM 可以利用题目文本和知识点描述弥补 ID 表征的语义不足；但其风险在于，语言生成的流畅性并不等价于教育有效性。因此，LLM-LPR 更需要显式的知识状态估计、检索约束、奖励对齐和可验证指标，否则容易退化为“看起来合理”的路径生成。

\begin{table}[t]
\centering
\caption{LPR 在不同教育应用系统中的典型形态}
\label{tab:lpr-applications}
\small
\begin{tabularx}{\textwidth}{p{2.5cm}p{2.8cm}p{2.8cm}X}
\toprule
应用系统 & 推荐动作 & 主要反馈 & 对 LPR 建模的影响 \\
\midrule
在线题库与智能辅导 & 下一题、下一知识点、短练习序列 & 对错、分数、用时、尝试次数 & 反馈较密集，适合闭环状态更新；关键在于把学习者状态、题目难度和知识结构结合起来\cite{Liu2019CSEAL,Zhang2024DLPR}。\\
MOOC 与课程平台 & 课程、章节、视频、学习单元 & 完成率、测验成绩、退课或跳过行为 & 动作粒度较粗、反馈稀疏、周期较长；更需要层级规划和路径级结构约束\cite{Zhang2019MOOC,Yu2025LIGHT,Chen2023SRC}。\\
考试准备与阶段学习计划 & 完整练习路径、复习计划、目标知识集合 & 阶段测验、目标掌握提升、覆盖率 & 需要一次性或半一次性生成可执行路径；难点是同时满足目标覆盖、难度平滑、预算限制和长期收益\cite{Chen2023SRC,Peng2026UNO,Cheng2026UniER}。\\
LLM 教师代理 & 路径、解释、提示、学习策略建议 & 学习收益指标、路径合理性、教师评价或代理奖励 & 输出从项目序列扩展到“推荐+解释+辅导”；必须用检索、知识状态和指标对齐约束语言生成\cite{Lv2025GenAL,Lv2025ReAL,Wang2026IBGRPO}。\\
\bottomrule
\end{tabularx}
\end{table}

总体而言，应用系统层面的差异可以解释许多方法设计选择：题库系统强调学生状态估计和题目难度，课程平台强调长期规划和层级动作空间，完整学习计划强调集合到序列的全局排序，LLM 教师代理则强调语义理解和教育目标对齐。把这些应用形态补入综述，有助于避免把 LPR 简化为单一算法问题，也能为后文讨论 KT、学生模拟器和 RL 路径规划提供更清晰的现实动机。

\section{现有 LPR 论文引言中的问题谱系}

从代表性 LPR 论文的引言可以看出，该领域的研究推动并不是简单地从一个模型替换为另一个模型，而是围绕“学习路径为什么难以推荐”这一核心问题不断细化。早期工作主要强调在线教育中的资源过载：学习者面对大量学习对象时，需要系统根据其目标、背景和偏好筛选合适资源。但 CSEAL 之后，问题重心逐渐从“推荐哪些资源”转向“如何在学习者状态动态变化的过程中，按教育逻辑安排资源顺序”。CSEAL 的引言明确指出认知结构包含两个互补部分：一是学习者对学习项目的知识掌握水平，二是学习项目之间的知识结构，例如先修关系。只利用前者的方法可能无法处理知识依赖，只利用后者的方法又无法体现不同学习者的能力差异。因此，它提出的问题不是一般意义上的个性化推荐，而是如何同时利用不可直接观测且持续演化的知识水平，以及相对稳定但复杂的知识结构\cite{Liu2019CSEAL}。

GEHRL 的引言进一步把目标导向性放在中心位置。它观察到，真实学习者往往不是随意刷题，而是带着一个或多个学习目标进入系统；不同目标之间存在难度差异和依赖关系，直接把目标集合的一次性编码交给 RL 智能体，会让策略在巨大动作空间中盲目探索。该论文因此把问题拆成两个子问题：先规划子目标的实现顺序，再为当前子目标推荐学习项目\cite{Li2023GEHRL}。这种写法背后的逻辑是：如果目标之间有先修或难度层次，那么“先学哪个目标”本身就是推荐问题的一部分，不能被简化为一个静态输入特征。

SRC 和 LIGHT 的引言则分别从完整路径生成和序列优化角度强调路径级建模的困难。SRC 认为，逐步推荐可以依赖即时反馈修正决策，但很多场景需要一开始就给出完整学习路径，例如课程组织或阶段性学习计划；完整路径生成搜索空间更大、反馈更少，因此比下一步推荐更受约束\cite{Chen2023SRC}。LIGHT 在此基础上指出，已有方法虽然考虑了显式先修关系或学习者反馈，却没有充分协调显式知识结构和隐式协同关系，导致路径可能失去逻辑连贯性\cite{Yu2025LIGHT}。二者共同说明：LPR 不只是一个动作选择问题，也是一个序列结构优化问题。

DLPR 的引言把问题推进到更细粒度的习题层。它指出，即便知识点顺序合理，路径中的习题难度仍可能剧烈波动：过难会让学习者挫败，过易会浪费学习预算；同样地，对不同难度的知识点分配相同练习次数也会造成效率低下。DLPR 因此把学习路径比作“行走路径”，认为学习项目难度相当于地形起伏，路径推荐必须同时追求平滑性和效率\cite{Zhang2024DLPR}。这一问题设定改变了 LPR 的评价重点：路径不只是知识点序列，还包含每一步练习是否处于学习者能力附近的问题。

GEPKSD、UNO 和近期 LLM 相关工作继续暴露出训练范式和表征层面的不足。GEPKSD 注意到，多学习者场景下环境动态不稳定：同一条路径对不同知识状态的学习者可能产生完全不同的结果；直接用 KT 预测完整知识状态再输入 RL 策略，既困难又可能冗余，于是它把真实知识状态视作训练阶段的特权信息，通过蒸馏把有用信息转移到可部署的推荐器中\cite{Li2024GEPKSD}。UNO 则认为现有 RL-LPR 常依赖终局奖励，无法判断路径中每一道题的具体贡献，并且按匿名会话训练会损失长期个性化信息\cite{Peng2026UNO}。GenAL、ReAL 和 IB-GRPO 进一步指出，传统 ID 表征忽略题目文本中的语义信息，而直接使用 LLM 又可能与教育目标不对齐，需要检索、代理或指标化 RL 约束来连接语义理解与长期学习收益\cite{Lv2025GenAL,Lv2025ReAL,Wang2026IBGRPO}。

因此，从这些论文的引言可以抽象出一条清晰的问题链：\textbf{资源过载}要求个性化；\textbf{先修关系和知识结构}要求顺序合理；\textbf{学习者状态动态变化}要求闭环决策；\textbf{多目标和长路径}要求层级规划；\textbf{习题难度和遗忘}要求细粒度过程建模；\textbf{稀疏奖励和离线日志}要求更稳定的训练范式；\textbf{题目语义和教育目标对齐}则推动 LLM/RAG 进入 LPR。后续方法的差异，本质上是对这条问题链中不同薄弱环节的回应。

\section{问题定义与数学逻辑}

\subsection{基本形式化}

设知识点集合为 $\mathcal{K}=\{k_1,\ldots,k_m\}$，学习项目集合为 $\mathcal{I}=\{i_1,\ldots,i_n\}$。学习项目可以是习题、视频、课程片段或学习对象。知识结构通常表示为有向图
\begin{equation}
\mathcal{G}=(\mathcal{K},\mathcal{E}),
\end{equation}
其中边 $(k_a,k_b)\in\mathcal{E}$ 表示 $k_a$ 是 $k_b$ 的先修知识，或二者存在某类教学相关关系。学习者 $u$ 的历史交互可表示为
\begin{equation}
\mathcal{H}_u=\{(i_t,r_t,\tau_t)\}_{t=1}^{T},
\end{equation}
其中 $i_t$ 是第 $t$ 次学习项目，$r_t$ 是反馈（如正误、分数、完成度），$\tau_t$ 是时间戳。知识追踪模型通常根据 $\mathcal{H}_u$ 估计学习者对各知识点的掌握状态
\begin{equation}
\mathbf{s}_t=f_{\mathrm{KT}}(\mathcal{H}_{u,1:t})\in [0,1]^m.
\end{equation}
给定目标知识集合 $\mathcal{K}^{\star}\subseteq\mathcal{K}$ 和学习预算 $B$，LPR 的目标是生成路径
\begin{equation}
\mathcal{P}_u=(a_1,a_2,\ldots,a_B),\quad a_t\in\mathcal{I}\ \text{或}\ \mathcal{K},
\end{equation}
使学习后的目标掌握收益最大化。一个常见抽象目标为
\begin{equation}
\max_{\pi}\ \mathbb{E}_{\pi}\left[U(\mathbf{s}_{B},\mathcal{K}^{\star})-\lambda C(\mathcal{P}_u)-\mu V(\mathcal{P}_u,\mathcal{G})\right],
\end{equation}
其中 $U$ 表示学习收益，$C$ 表示路径成本或长度，$V$ 表示违反先修关系、难度跨度或重复冗余带来的惩罚。

这个形式化背后的数学逻辑较简单：推荐策略 $\pi$ 的每一次动作都会改变学习者状态，因此路径价值不是单步收益的简单相加，而是由状态转移累积形成。若用马尔可夫决策过程（MDP）表示，则有
\begin{equation}
\mathcal{M}=(\mathcal{S},\mathcal{A},P,R,\gamma),
\end{equation}
其中状态 $s_t$ 包含当前知识掌握、目标、剩余预算和时间信息；动作 $a_t$ 是推荐知识点或学习项目；转移 $P(s_{t+1}|s_t,a_t)$ 由学习行为和学生模型决定；奖励 $R(s_t,a_t)$ 则衡量即时或延迟学习收益。强化学习式 LPR 的难点在于：真实教育反馈往往稀疏、延迟且受个体差异影响，导致奖励设计和离线评估比一般推荐任务更困难。

\subsection{与知识追踪的关系}

KT 和 LPR 的关系可概括为“诊断--决策”闭环。KT 解决“学生现在会什么”的问题，LPR 解决“下一步应该学什么”的问题。若 KT 状态估计偏差较大，则 LPR 会在错误状态上进行规划。反过来，LPR 的推荐动作会决定后续 KT 能观察到什么数据，形成反馈回路。因此，一个严格的 LPR 系统通常需要同时考虑
\begin{equation}
\mathbf{s}_{t+1}=f_{\mathrm{KT}}(\mathbf{s}_t,a_t,r_t,\Delta t),
\end{equation}
以及
\begin{equation}
a_t\sim \pi_{\theta}(a|\mathbf{s}_t,\mathcal{K}^{\star},B-t).
\end{equation}
前者是学习者状态演化，后者是路径推荐策略。二者共同决定最终学习效果。

\begin{figure}[t]
\centering
\begin{tikzpicture}[
  node distance=1.1cm,
  box/.style={draw,rounded corners,align=center,minimum width=2.7cm,minimum height=0.8cm,fill=gray!8},
  main/.style={draw,rounded corners,align=center,minimum width=3.0cm,minimum height=0.9cm,fill=blue!7},
  arr/.style={-{Latex[length=2mm]},thick}
]
\node[main] (hist) {学习历史\\$\mathcal{H}_u$};
\node[main,right=of hist] (kt) {知识追踪\\$\mathbf{s}_t=f_{\mathrm{KT}}(\mathcal{H}_u)$};
\node[main,right=of kt] (policy) {路径策略\\$\pi(a_t|\mathbf{s}_t)$};
\node[box,below=of policy] (item) {推荐项目\\知识点/习题};
\node[box,below=of kt] (feedback) {作答反馈\\$r_t,\Delta t$};
\node[box,below=of hist] (goal) {学习目标\\$\mathcal{K}^{\star},B$};
\draw[arr] (hist) -- (kt);
\draw[arr] (kt) -- (policy);
\draw[arr] (goal) -- (policy);
\draw[arr] (policy) -- (item);
\draw[arr] (item) -- (feedback);
\draw[arr] (feedback) -- (kt);
\draw[arr,dashed] (feedback) -- (hist);
\end{tikzpicture}
\caption{学习路径推荐与知识追踪的闭环关系。KT 提供状态估计，LPR 根据状态和目标选择动作，推荐动作产生新反馈并再次更新 KT 状态。}
\label{fig:kt-lpr-loop}
\end{figure}

\section{CSEAL 之前：传统算法与早期路径推荐}

CSEAL 之前的 LPR 研究并非没有序列意识，但多数方法更接近“结构约束下的路径搜索”或“学习对象推荐”，而不是显式的长期策略优化。其共同特点是把知识结构、学习对象属性、学习者偏好或历史行为作为约束，再通过规则、启发式搜索或监督模型生成路径。

\subsection{本体、知识图与规则驱动方法}

本体和语义网方法较早被用于上下文感知的 e-learning 推荐。Yu 等通过本体表示学习对象、上下文和用户偏好，使系统能够根据语义关系进行资源匹配\cite{Yu2007Ontology}。这类方法的优势在于可解释性强：推荐理由可直接追溯到本体关系、先修约束或学习对象属性。但其局限也明显，即知识结构需要人工构建和维护，面对大规模、快速变化的课程资源时成本较高。

知识图或知识地图方法进一步把学习对象组织为图结构。Zhu 等提出基于知识地图的多约束路径推荐，把学习目标、先修关系、学习者能力和资源约束纳入路径规划\cite{Zhu2018KnowledgeMap}。这类方法可以看作约束优化问题：在图上寻找满足条件的路径，使路径代价最小或学习收益最大。其数学核心接近
\begin{equation}
\min_{\mathcal{P}}\ \sum_{a_t\in\mathcal{P}}c(a_t)
\quad \mathrm{s.t.}\quad
\mathcal{P}\ \text{满足先修、难度、时间等约束}.
\end{equation}
相较于后来的 RL 方法，它们通常缺少学习者状态的动态闭环更新。

\subsection{协同过滤、序列推荐与深度序列模型}

另一条早期路线来自推荐系统。协同过滤和 KNN 类方法通过相似学习者或相似资源进行推荐，具有实现简单和可扩展的优势，但难以表达先修依赖。随着深度序列模型发展，LSTM、GRU4Rec 等模型被用于建模学习行为序列。Zhou 等提出基于 LSTM 的个性化学习全路径推荐模型，通过历史学习序列预测后续路径\cite{Zhou2018LSTM}；面向 MOOC 的课程推荐也开始尝试层级强化学习，以处理课程级长期选择问题\cite{Zhang2019MOOC}。这类方法增强了序列建模能力，但其优化目标通常是复现历史行为、预测下一项或优化课程选择，而不是直接最大化细粒度知识掌握收益。

\subsection{组合优化与学习对象路径规划}

传统优化方法把 LPR 视为组合搜索问题，常见策略包括遗传算法、蚁群优化、多约束路径搜索和辅助学习对象选择。Nabizadeh 等的辅助学习对象推荐强调根据学习者状态补充必要资源\cite{Nabizadeh2020Auxiliary}；综述工作则系统总结了学习路径个性化中的图搜索、规则、群智能和推荐算法\cite{Nabizadeh2020Survey}。这些方法在可解释性和约束可控性上仍有价值，尤其适合学习目标明确、知识图可靠且数据规模有限的场景。

总体而言，CSEAL 之前的传统方法为 LPR 提供了三个基础：知识结构、路径约束和学习者个性化特征。但它们通常没有把“推荐--反馈--再推荐”建模为在线序列决策，也缺少与 KT 状态估计的紧密结合。

\section{知识追踪、学生模拟器与 RL-LPR 的状态基础}

在 CSEAL 之后的强化学习式 LPR 中，KT 不应只被视为一个外围预测模块。更准确地说，RL-LPR 通常需要一个“学生模型”来回答三个问题：第一，学习者当前处于什么知识状态；第二，推荐某个知识点或习题后，学习者状态如何变化；第三，系统应如何把这种状态变化转化为奖励。KT、IRT 以及后续的 KES 模型分别从不同角度承担这些功能。因此，把 KT 相关内容放在强化学习式 LPR 之前讨论更合理，因为后续 CSEAL、GEHRL、DLPR、UNO 以及更细粒度的逐题闭环式路径推荐思路，都依赖某种学生状态估计或模拟机制。

\subsection{从 IRT/KSS 到 KT/KES：学生模型的角色变化}

CSEAL 中使用的 KSS（Knowledge State Simulator）并不是典型的深度 KT 模型，而更接近基于 IRT 思想的学生知识状态模拟。IRT 的基本逻辑是用学习者能力和题目参数解释作答结果。以最简单的一参数或二参数 IRT 为例，学习者 $u$ 正确回答题目 $i$ 的概率可写成
\begin{equation}
P(r_{ui}=1)=\sigma(a_i(\theta_u-b_i)),
\end{equation}
其中 $\theta_u$ 表示学习者能力，$b_i$ 表示题目难度，$a_i$ 表示区分度，$\sigma(\cdot)$ 是 sigmoid 函数。IRT 的优势是可解释，能把能力、难度等教育测量变量显式建模；不足是通常更偏静态，难以充分表达长期序列交互中的知识迁移、遗忘和多知识点联动。CSEAL 采用 KSS 的意义在于为 RL 环境提供可交互反馈，使 Actor-Critic 策略可以在模拟环境中评估推荐动作，而不是直接在真实学生上高成本探索。

KT 模型则更强调时间序列状态更新。DKT 用 RNN 从历史作答序列中学习隐状态\cite{Piech2015DKT}，GKT 将知识点关系图引入状态传播\cite{Nakagawa2019GKT}，CKT 和 RKT 借助注意力、上下文与题目关系建模遗忘和关联影响\cite{Ghosh2020CKT,Pandey2020RKT}，DIMKT 进一步显式考虑题目难度对知识状态变化的影响\cite{Shen2022DIMKT}。近期 DLKT 综述和过程一致性 KT 研究也说明，KT 正在从单纯追求下一题预测准确率，转向更关注学生知识状态演化是否符合学习过程\cite{Liu2025DLKTReview,Shen2024LPKT}。如果说 IRT/KSS 更像是“基于能力参数的学生反应模型”，那么 KT/KES 更像是“基于学习历史的状态转移模型”。在 RL-LPR 中，后者的重要性在于它可以近似给出
\begin{equation}
\mathbf{s}_{t+1}=f_{\mathrm{student}}(\mathbf{s}_t,a_t,r_t,\Delta t),
\end{equation}
从而让推荐策略具有可训练的环境反馈。

近年来许多 LPR 工作使用 KES（Knowledge Evolution Simulation 或 Knowledge Estimation Simulator）类模型作为学生模拟器。KES 的任务不是单纯预测下一题正确率，而是在模拟推荐交互时给出作答反馈和状态更新，使 RL 智能体能够在离线环境中训练。例如 DLPR 使用 KES-DIMKT 作为学生模拟器来支持难度感知路径规划\cite{Zhang2024DLPR}。这种做法的本质是用 KT/KES 近似真实学习环境：策略选择动作，模拟器返回反馈，状态估计器更新学生状态，奖励函数再根据目标掌握变化评价动作。

\subsection{KT 在 LPR 中的三种用法}

现有 LPR 对 KT 的使用大致可分为三类。第一类是\textbf{状态估计器}：KT 根据学生历史输出知识点掌握向量 $\mathbf{s}_t$，作为 RL 策略或序列生成模型的输入。CSEAL 的知识水平估计、GEHRL 的学习者状态建模、DLPR 的难度感知状态估计都属于这一类。第二类是\textbf{学生模拟器}：KT/KES 模拟学生对推荐项目的反馈，使策略可以在离线环境中训练，避免真实在线探索的成本。第三类是\textbf{奖励评估器}：KT 估计推荐前后目标知识点掌握度的变化，用于构造学习收益奖励。

这三类用法之间有明显差异。作为状态估计器时，KT 主要影响策略输入；作为学生模拟器时，KT 决定训练环境是否可信；作为奖励评估器时，KT 直接决定策略优化方向。如果同一个 KT 模型同时承担这三种角色，就可能出现闭环偏差：策略在一个由 KT 定义的环境中训练，又用同一个 KT 的输出评价自己，最终得到的高收益未必等价于真实学生收益。因此，严格的 LPR 研究需要区分“真实学生状态”“模拟器状态”和“推荐器可观测状态”。GEPKSD 将特权知识状态与部署时可获得的练习日志区分开来，正是对这一问题的回应\cite{Li2024GEPKSD}。

\subsection{KT 指标与 LPR 目标之间的错位}

KT 常用 AUC、ACC 等指标评价下一题预测，而 LPR 更关注策略作用后的累计学习收益。二者有关联，但不是同一目标。一个 KT 模型可能具有较高 AUC，却不一定能产生对路径决策最有用的状态表示；相反，一个状态表示即使不精确恢复每个知识点掌握概率，只要能区分“下一步推荐什么最有效”，也可能足够支持 LPR。这正是 GEPKSD 质疑“完整知识状态预测是否必要”的原因。

从数学上看，KT 优化的是
\begin{equation}
\min_{\phi}\sum_t \ell(\hat{r}_{t+1},r_{t+1}),
\end{equation}
而 LPR 优化的是
\begin{equation}
\max_{\pi}\mathbb{E}_{\pi}\left[\sum_{t=1}^{B}\gamma^{t-1}R(\mathbf{s}_t,a_t)\right].
\end{equation}
前者关注预测误差，后者关注动作选择后的累计收益。如果奖励 $R$ 本身由 KT 状态变化构造，那么 KT 的误差会进入策略训练和策略评估两个环节。因此，未来更合理的方向不是简单追求更高 KT AUC，而是研究哪些状态维度对路径推荐决策最关键，例如目标知识点掌握度、先修薄弱点、题目级正确概率、遗忘风险、难度适配和剩余预算。

\subsection{逐题闭环式 LPR 对学生模型使用方式的启示}

上述研究也引出一个更细粒度的学生模型使用思路：未来 RL-LPR 可以把状态估计、学生模拟和局部习题评估拆成不同职责，而不是让同一个黑盒学生模型同时承担所有功能。例如，可以使用一个状态估计模块为知识点规划器和局部习题执行器提供 $(h_t,p_t)$ 等状态，同时使用独立学生模拟器提供作答反馈，避免把状态估计器和训练环境完全混同。该思路强调：细粒度 LPR 不应只在知识点层做规划，也不应只依赖一个黑盒模拟器给出终局收益，而应在每一道习题后更新状态、重新规划知识点，并用局部候选题的一步前瞻收益监督习题选择。

在这一视角下，KT/KES 不只是“给 RL 一个环境”的工具，而是连接知识点规划和习题执行的反馈通道。知识点规划器需要知道目标集合中哪些知识点仍未达标、哪些已达标但需要巩固；局部执行器则需要知道当前学生对候选习题和相关知识点的掌握状态。一个值得进一步研究的方向是把这两个粒度统一到逐题闭环中：每完成一道题，学生状态被更新，高层规划器可以继续、切换或回访知识点，局部执行器再围绕新知识点选择具体习题。这种设计体现了一种更细粒度的观点：有效 LPR 的核心不只是推荐一串知识点，而是让“知识点规划--习题执行--状态更新”形成可调整的循环。

\section{CSEAL 之后：强化学习与深度路径规划}

\begin{figure}[t]
\centering
\begin{tikzpicture}[
  font=\small,
  stage/.style={
    draw,
    rounded corners,
    align=left,
    text width=0.84\textwidth,
    minimum height=0.95cm,
    inner xsep=0.35cm,
    inner ysep=0.12cm,
    fill=blue!5
  },
  newstage/.style={
    stage,
    fill=orange!9
  },
  arr/.style={-{Latex[length=2mm]},thick}
]
\node[stage] (m0) {\textbf{2007--2018：传统路径规划} \quad 本体、知识图、多约束优化和序列推荐主要解决资源组织与可解释路径问题。};
\node[stage,below=0.32cm of m0] (m1) {\textbf{2019：CSEAL} \quad 将认知结构、MDP 和 Actor-Critic 结合，推动 LPR 从静态路径搜索转向动态序列决策。};
\node[stage,below=0.32cm of m1] (m2) {\textbf{2023：GEHRL / SRC} \quad 分别从目标导向层级规划和集合到序列生成角度处理多目标路径推荐。};
\node[stage,below=0.32cm of m2] (m3) {\textbf{2024：DLPR / GEPKSD} \quad 引入难度感知、层级图建模和特权知识状态蒸馏，关注细粒度练习与多学习者动态。};
\node[newstage,below=0.32cm of m3] (m4) {\textbf{2025：LIGHT / KnowLP / GenAL} \quad 融合显式先修、隐式协同、GraphRAG 自动建图和 LLM 语义理解。};
\node[newstage,below=0.32cm of m4] (m5) {\textbf{2026：UNO / IB-GRPO / UniER} \quad 面向离线训练、过程奖励、教育目标对齐和统一基准评测。};
\draw[arr] (m0) -- (m1);
\draw[arr] (m1) -- (m2);
\draw[arr] (m2) -- (m3);
\draw[arr] (m3) -- (m4);
\draw[arr] (m4) -- (m5);
\end{tikzpicture}
\caption{学习路径推荐技术演进示意。CSEAL 之前主要是结构和优化方法；CSEAL 之后，强化学习、图增强、难度感知、离线训练和大模型增强逐步成为主线。}
\label{fig:evolution}
\end{figure}

\subsection{CSEAL：认知结构与 Actor-Critic 的结合}

CSEAL 是 LPR 领域的重要分水岭。它首先明确指出，早期自适应学习方法常把认知结构拆开使用：有的方法只估计学习者知识水平，却不处理学习项目间的先修依赖；有的方法只利用知识结构，却忽略不同学习者的当前能力和学习节奏。前者容易推荐出违背知识依赖的学习项目，后者则可能给不同水平学习者安排相同路径。CSEAL 认为，有效路径推荐必须同时利用两类认知信息：动态变化、不可直接观测的知识水平，以及相对稳定、体现先修或相关关系的知识结构\cite{Liu2019CSEAL}。

围绕这一问题，CSEAL 将自适应学习建模为 MDP。状态用于刻画学习者当前知识水平，动作对应下一步推荐的学习项目，奖励来自学习效果提升。由于学习者真实掌握水平不能直接观测，CSEAL 需要根据历史测试或练习反馈估计状态；由于动作空间中存在大量教学上不合理的项目，它又引入知识结构导航来限制或引导策略选择。其方法设计可以理解为三层逻辑：
\begin{itemize}
  \item 用学生历史交互估计认知状态，解决“学习者当前会什么”不可直接观测的问题；
  \item 用知识结构限制或引导推荐动作，避免策略只从统计相关性出发而推荐不符合先修逻辑的项目；
  \item 用 Actor-Critic 优化长期学习收益，而不是只预测下一道题或下一项资源。
\end{itemize}
为什么这种设计能够缓解早期方法的问题？关键在于它把“学生状态”和“知识结构”放入同一个决策闭环中。知识状态使策略能够区分不同学习者，知识结构使策略不会完全退化为黑盒推荐，RL 则使系统可以评价一段路径的长期收益。它的局限也同样清晰：状态估计依赖学生模拟或历史反馈质量，认知导航依赖已有知识结构，奖励仍更多关注可观测学习效果。因此，后续工作大多是在这三个薄弱点上继续推进：更好的目标规划、更细的习题难度、更稳定的训练和更自动化的知识结构构建。

\subsection{层级强化学习：从知识点到习题}

GEHRL 的出发点是 CSEAL 之后仍未充分解决的“目标使用不足”问题。真实学习者往往带着一个或多个目标进入系统，例如想掌握“微分”或同时掌握“泰勒公式、随机变量和微积分”。这些目标之间并不是平行的：有的目标更简单，有的目标依赖其他知识点，有的目标需要先通过若干子目标铺垫。GEHRL 认为，若只把目标集合编码为 one-hot 向量输入策略网络，智能体需要在大动作空间中自行探索目标实现顺序，容易出现两类问题：一是缺少目标规划，不知道先实现哪个目标更合适；二是实现某个目标时可能推荐大量与目标无关的学习项目，降低学习效率\cite{Li2023GEHRL}。

为此，GEHRL 将目标导向 LPR 拆成“子目标选择”和“子目标达成”两个层次。高层智能体负责从学习项目中选择当前子目标：它既可以安排多个目标的实现顺序，也可以把困难目标拆成更简单的先修子目标。低层智能体则围绕高层给出的子目标，逐步推荐具体学习项目。图结构在其中的作用不是装饰性的特征，而是帮助高层理解目标之间的依赖，帮助低层缩小与当前子目标相关的搜索空间。

这种分层的数学动机是降低动作空间复杂度。如果直接在习题集合 $\mathcal{I}$ 上决策，动作空间大小为 $|\mathcal{I}|$；若先选择知识点再选择习题，则高层空间为 $|\mathcal{K}|$，低层只需在与当前知识点相关的候选集合 $\mathcal{I}(k)$ 中搜索。层级结构近似为
\begin{equation}
\pi(a_t|s_t)=\pi_{\mathrm{low}}(i_t|s_t,k_t)\pi_{\mathrm{high}}(k_t|s_t).
\end{equation}
该分解与教师教学中的“先确定学什么，再选择怎样练”一致。它认为层级规划能够解决 CSEAL 式扁平策略对多目标结构利用不足的问题：高层先决定目标顺序，低层再保证当前路径尽量与目标相关。不过，这种框架也引入新的复杂性：高层子目标如果选错，低层即使推荐得很好也可能偏离最终目标；低层反馈如何稳定传递给高层，也会影响整体训练效率。

DLPR 则发现，即使层级规划解决了“学什么”的问题，仍未充分解决“练什么难度”的问题。其引言把学习过程类比为真实行走：如果一条路地形忽高忽低，行走者会感到吃力；类似地，如果学习路径中题目难度突然升高或突然降低，学习者可能产生挫败或无聊。另一方面，不同难度知识点如果被分配相同练习次数，也会造成学习资源浪费：太难的项目练几次就切换可能无法掌握，太简单的项目反复练习又会降低效率\cite{Zhang2024DLPR}。

DLPR 因此把学习项目拆为学习项目和练习项目两个层级，并构建层级图来表达二者关系，再通过难度驱动的层级强化学习进行决策。高层 L-Agent 选择接下来要练习的学习项目，低层 P-Agent 在该学习项目相关的练习项目中选择具体题目。与 GEHRL 相比，DLPR 的关键不是单纯分层，而是在状态估计、代理通信和奖励约束中显式注入难度信息，使高层和低层都知道当前路径是否过陡、过平或练习不足。其核心思想可以抽象为
\begin{equation}
R_t=R_{\mathrm{gain}}(s_t,i_t)-\eta |d_{i_t}-\alpha_t|.
\end{equation}
其中 $R_{\mathrm{gain}}$ 表示学习收益，$d_{i_t}$ 表示练习难度，$\alpha_t$ 表示当前能力估计。这个式子并不一定是原文的完整奖励函数，但体现了 DLPR 的数学逻辑：学习收益需要受到难度适配项约束。难度项不是附属特征，而是决定学习者是否能在“最近发展区”内获得有效练习的重要变量。

\subsection{集合到序列与知识拓扑优化}

SRC 的问题意识不同于 CSEAL 和 GEHRL。它首先区分两类 LPR 场景：一类是每一步根据学习者反馈实时推荐下一项，另一类是在学习开始时一次性给出完整路径。前者可以利用即时反馈修正策略，后者则更适合课程规划、阶段性学习安排等需要提前知道完整路线的场景。但完整路径生成更困难，因为路径搜索空间更大，可利用的中间反馈更少，并且目标概念集合本身没有天然顺序\cite{Chen2023SRC}。

因此，SRC 将任务建模为 set-to-sequence：输入是一个无序的目标概念集合，输出是有序学习路径。这样建模的原因在于，目标集合的排列顺序不应影响模型理解，但输出路径必须有顺序。SRC 通过概念感知编码器捕获目标概念之间的关联，再由解码器逐步生成路径；同时引入辅助 KT 模块来估计路径对学习效果的影响，并用策略梯度优化排序目标。它试图解决的问题不是“下一步学什么”，而是“给定一组目标，怎样生成一条整体上合理的学习序列”。与逐步 RL 相比，SRC 的优势在于路径级训练和并行建模更方便，弱点则是路径执行中的实时反馈和中途重规划能力相对不足。

LIGHT 进一步指出，无论是逐步推荐还是完整路径生成，现有方法往往没有充分协调两类关系：显式关系，例如先修结构；隐式关系，例如从学习者历史交互中体现出来的协同学习信号。只使用显式先修关系，可能忽略学习者数据中反复出现的共学模式；只使用隐式协同关系，又可能生成缺乏教学逻辑的路径。LIGHT 因此把 LPR 重新放到知识拓扑感知的序列优化框架中\cite{Yu2025LIGHT}。

具体来说，LIGHT 先构建复合概念图，将显式先修关系和隐式协同关系放在两个互补视角中；然后通过互补对比融合学习更稳健的概念表示，使模型既保留教学结构，又吸收学习行为中的统计信号。接着，它通过结构语义聚类和候选路径采样降低路径搜索空间，再用双向感知的路径优化网络从序列角度评价候选路径。它认为这样可以同时解决两个问题：一是学习元素本身不是孤立节点，需要在图拓扑中学习表示；二是完整路径空间随路径长度指数增长，必须通过拓扑采样和序列优化提高搜索效率。该类方法说明，路径合理性不能只由单一先修边决定。两个知识点之间可能存在先修、相似、共现、难度连续或学习迁移等多种关系；多关系图更符合真实学习环境。

\subsection{知识状态蒸馏、离线训练与奖励设计}

GEPKSD 关注的是 RL-LPR 在多学习者环境中的训练不稳定问题。它指出，同一条学习路径对不同学习者可能产生不同结果，因为学习者的初始知识水平和学习动态不同；这使得 RL 智能体面对的是一个动态不稳定的环境。已有方法通常用 KT 模型从练习日志中预测细粒度知识状态，再把该状态输入推荐策略。但 GEPKSD 对这种做法提出质疑：路径推荐确实需要反映学习者知识状态，却未必需要准确恢复每个概念的完整掌握概率；如果这些概率最终还要被另一个神经网络压缩成隐向量，那么直接学习其中对推荐有用的信息可能更高效\cite{Li2024GEPKSD}。

因此，GEPKSD 将真实或模拟器提供的知识状态视为训练阶段的特权信息。训练时，模型可以访问更强的知识状态编码，并学习它如何帮助推荐；推理时，真实知识状态不可得，于是用知识状态适配器从普通练习日志中提取类似的隐信息，替代特权状态编码。这个思路本质上是在解决部分可观测 MDP 问题：真实学生能力 $\theta_t$ 不可直接观测，系统只能看到交互序列和模型估计状态。它不要求部署阶段精确预测 $\theta_t$，而是要求部署模型复现“特权状态对推荐有用的那部分信息”。

UNO 则把问题转向训练范式和奖励密度。它认为现有 RL-LPR 常依赖终局奖励，即完整路径结束后才根据最终测试或掌握提升评价整条序列。这会产生典型的信用分配问题：如果整条路径效果不好，其中某些真正有效的题目也可能被误判为无效；如果整条路径效果好，也难以知道哪些步骤贡献最大。另一个问题是匿名会话训练：把长期学习记录切成匿名 session，会损失学生跨会话的个性化信息，导致策略更像是在学习平均路径，而不是长期适应某个学习者\cite{Peng2026UNO}。

UNO 因此提出统一离线训练范式，用更密集的过程奖励刻画每一步对学习目标的贡献，并从长期历史中构造更个性化的训练样本。在真实教育场景中，在线探索成本很高，错误推荐可能伤害学习体验，因此离线 RL 和保守策略优化具有实际意义。数学上，离线 LPR 要在固定日志数据 $\mathcal{D}$ 上学习策略 $\pi$，但不能让新策略偏离数据支持范围太远，否则价值估计会失真。该问题与一般离线 RL 中的分布外动作风险相同。

此外，遗忘因素与知识图谱结合的路径推荐也在近期出现。Fan 等将遗忘因素、知识图谱感知和 RL 结合，强调推荐不应只关注新知识推进，也应考虑已学知识的保持与回访\cite{Fan2025ForgettingKG}。这与 KT 中的时间衰减和遗忘建模形成呼应，说明长期学习路径需要同时处理“学新”和“防忘”。从问题逻辑上看，遗忘建模补足了许多路径推荐默认假设中的漏洞：学习者掌握某知识点并不意味着该状态永久保持，路径规划必须考虑复习时机和知识保持收益。

\subsection{从路径生成到知识点--习题循环规划}

在上述工作基础上，可以进一步引出一种知识点--习题循环规划思路：现有 LPR 虽然逐渐从知识点级规划走向习题级推荐，但“知识点规划”和“习题执行”之间的反馈闭环仍不够充分。知识点级方法可以降低动作空间并保持路径结构，但难以区分同一知识点下不同习题的学习作用；习题级方法虽然能输出具体练习，却可能把某个知识点的学习视为相对独立的局部过程，未能充分利用每道题反馈及时调整后续知识点规划。换言之，细粒度 LPR 的关键不只是把推荐对象从知识点下沉到习题，而是让每一次习题反馈都能重新影响知识点层面的决策。

这种潜在设计可以分为三部分。第一，状态估计模块同时估计学生在知识点级和题目级的状态，使规划器能够看到目标知识点掌握度和具体题目作答概率。第二，高层规划器在每道题之后重新规划知识点，而不是在一个固定子会话内保持高层目标不变；因此它可以根据最新反馈选择继续练习、切换知识点或回访旧知识点。第三，局部习题执行器不是简单选择难度匹配或历史相似的题目，而是用一步前瞻收益构造监督信号，让匹配模型学习“在当前状态下哪道候选题最有助于当前知识点目标”。

这种思路对奖励设计也提出了一个值得纳入综述的观点。UNO 等工作已经指出终局奖励过于稀疏，但若只使用离散达成奖励，即只有当知识点跨越掌握阈值时才给正反馈，仍会忽略未跨阈值前的渐进提升和达标后的巩固收益。因此，一个自然方向是引入离散--连续混合过程奖励：离散部分刻画目标是否达标，连续部分刻画目标集合整体掌握度的细微变化。其逻辑可概括为
\begin{equation}
R_t^{\mathrm{hybrid}}
=
\Delta E_t^{N}
+
\Delta E_t^{S}
-
\eta,
\end{equation}
其中 $\Delta E_t^{N}$ 表示离散达成收益，$\Delta E_t^{S}$ 表示连续掌握收益，$\eta$ 表示步级成本。这个奖励思想体现了一个基本看法：LPR 不应只在路径终点评价成败，也不应只在跨越阈值时承认学习有效；只要学生在目标知识上发生可观测的连续改善，就应当成为策略学习的信号。

从方法谱系上看，知识点--习题循环规划可以被理解为对 CSEAL、GEHRL、DLPR 和 UNO 所暴露问题的一种综合回应：它继承 CSEAL 的动态闭环思想，吸收 GEHRL/DLPR 的层级决策视角，回应 UNO 对稀疏奖励的批评，同时把局部习题选择从 RL 策略动作中适当拆出，交给监督式执行器或匹配模型承担。这样做的好处是降低 RL 在巨大题目空间中的探索压力，同时保留习题层推荐的细粒度收益。不过，这类方向仍依赖学生模拟器和 KT/KES 状态估计，其真实在线效果、跨数据集泛化和解释能力仍需要进一步验证。

\section{大模型、RAG 与知识图谱增强}

\subsection{GraphRAG：从人工知识图到自动知识结构构建}

近年的趋势之一是利用大语言模型和检索增强技术构建、解释或执行学习路径推荐。KnowLP 的问题设定是：先修关系对路径推荐很重要，但高质量先修图通常依赖专家标注，成本高、覆盖有限；即便存在先修图，也可能稀疏、有噪声，并且过于刚性。更重要的是，只沿先修关系推荐时，如果学习者在某个概念上受阻，系统缺少帮助其区分相似概念的机制。KnowLP 因此引入 GraphRAG 自动诱导双知识结构图，包括先修关系图和相似关系图\cite{Cheng2025KnowLP}。

KnowLP 的核心逻辑不是简单让 LLM 生成一张知识图，而是把“先修学习”和“辨析学习”结合起来。先修图用于维持学习路径的阶段性和方向性，相似图用于在学习者遇到混淆概念时引入对比支架，帮助其区分相近知识点。其 EDU-GraphRAG 模块用于从目标数据中构建任务相关知识结构，并通过优化反馈降低 LLM 幻觉；随后，强化学习模块中的先修代理和相似代理分别决定是否遵循先修依赖、是否引入相似概念进行辨析。这样做的原因是：如果只依赖先修边，路径可能在某个难点上阻塞；如果加入相似概念，学习者可以通过对比学习恢复路径推进。

这一类工作对传统 LPR 的改变在于，知识结构不再只是论文预先给定的静态输入，而可能成为模型自动构建和动态修正的对象。其优势是降低人工标注成本，并能为不同课程或数据集生成更贴近局部任务的知识关系；风险则是自动生成的边未必具有教育真实性。一个合理的评估方式不应只看生成图的密度或覆盖率，而应检查生成关系是否能解释学习者表现。例如，若 $k_a$ 被判定为 $k_b$ 的先修，则未掌握 $k_a$ 的学生在 $k_b$ 上的成功率应显著受限；若两个知识点被判定为相似或易混淆，则在错误模式、题目文本或迁移表现中应能找到相应证据。

\subsection{LLM 代理：从 ID 表征到语义驱动的教学决策}

GenAL 和 ReAL 进一步指出，传统 LPR 大多把题目、知识点和学习者表示为 ID 或稀疏标签，这会丢失题目文本中的细粒度语义。例如两道题都标注为“线性函数”，但一道可能只考察基础概念，另一道可能同时涉及图像、距离公式和综合推理；ID 表征很难区分这些认知层次。GenAL 因此把生成式代理引入自适应学习，希望利用 LLM 的语义理解能力解析题目文本、学习者状态和知识目标\cite{Lv2025GenAL}。ReAL 则进一步强调检索增强代理的作用：单靠 LLM 内部知识可能不稳定，结合检索到的题目、知识点和教师式反馈，可以让代理更接近真实教师的决策过程\cite{Lv2025ReAL}。

LLM 代理在 LPR 中可能承担三类任务。第一是\textbf{语义解析}：从题目文本中识别隐含知识点、难度层级、解题步骤和认知要求。第二是\textbf{候选生成}：根据学习目标和当前状态生成可供传统策略模型筛选的候选路径或候选习题。第三是\textbf{解释生成}：把“为什么推荐该知识点或该题”的理由转换成教师和学生可理解的自然语言。与纯 RL 策略相比，LLM 的优势在于语义泛化和冷启动；与纯 LLM 推荐相比，RL/KT 的优势在于可以用学生状态和学习收益约束决策。因此，更稳妥的方向可能不是让 LLM 独立决定完整路径，而是让 LLM 负责语义增强，路径选择仍由可评估的策略模块完成。

\subsection{教育目标对齐：从语言合理到学习有效}

IB-GRPO 则从大模型路径推荐与教育目标对齐出发。它认为，LLM 虽然具有语义理解能力，但预训练目标是 next-token prediction，并不天然等价于长期学习收益最大化。直接用提示词生成路径，可能得到语言上合理但教学上短视的推荐；同时，真实 LPR 还要同时满足最近发展区、路径长度、难度调度、多样性和最终学习效果等多个目标。IB-GRPO 因此使用指标驱动的 group relative policy optimization，把学习效果、难度调度、长度可控和轨迹多样性等指标显式纳入优化\cite{Wang2026IBGRPO}。这些工作说明，LLM 在 LPR 中不只是文本生成器，也可能承担知识结构构建、路径解释、教师代理和策略优化器等角色。

但 LLM-LPR 仍应谨慎看待。第一，教育路径推荐是高后果序列决策，错误推荐会累积影响学习效果。第二，LLM 的自然语言解释可能看似合理但事实不可靠。第三，若缺少真实学习反馈，LLM 生成路径的教育有效性难以验证。因此，大模型增强 LPR 的关键不在于“能生成路径”，而在于能否被可验证的教育目标约束。更稳妥的方向是将 LLM 放在可审计的闭环中：让它负责语义解析、候选解释或知识结构补全，而路径选择仍由可评估的学习收益、难度约束和学生状态更新共同约束。

从逐题闭环式细粒度 LPR 的角度看，LLM/RAG 可以与知识点--习题循环规划形成互补关系。逐题闭环关注知识点规划与习题执行之间的反馈耦合，而 LLM 可以进一步增强习题文本理解和推荐解释。例如，局部执行器在构造候选题时，可以引入 LLM 提取的题目语义、认知层级或解题步骤；在输出推荐时，也可以根据当前知识状态、目标知识点和候选题特征生成自然语言解释。但这类融合仍应受 KT/KES 状态更新和学习收益指标约束，否则容易退化为“语义上看似合理、教育上未必有效”的路径生成。

\section{常用数据集与基准选择}

虽然本文不是实验论文，但数据集选择仍然需要单独讨论。原因在于，LPR 的研究结论往往高度依赖数据集能够观测到什么。一个教育交互数据集可以抽象为
\begin{equation}
\mathcal{D}=\{(u_t,i_t,c_t,y_t,\tau_t,\mathbf{x}_{i_t})\}_{t=1}^{N},
\end{equation}
其中 $u_t$ 是学习者，$i_t$ 是题目、课程或学习资源，$c_t$ 是知识点或技能标签，$y_t$ 是答题结果、分数或完成行为，$\tau_t$ 是时间信息，$\mathbf{x}_{i_t}$ 表示题目文本、难度、先修关系等辅助信息。若数据集中只有 $(u_t,i_t,y_t)$，模型最多能较稳妥地学习“谁可能答对什么”；若还包含知识点、时间、题目文本和关系结构，才更可能支撑知识追踪、难度调度、语义增强和路径级推荐。因此，数据集并不是实验部分的附属物，而是决定 LPR 可建模变量边界的前提。

对 LPR 来说，数据集至少影响四个关键问题。第一，\textbf{动作粒度}：ASSISTments、Junyi、KDD Cup 等数据集多是习题级交互，因此适合研究下一题推荐、练习序列和路径级 exercise recommendation；MOOC 数据则更适合课程或章节级规划。第二，\textbf{状态可观测性}：若数据集包含知识点标签和较密集的答题反馈，KT 或 KSS/KES 模拟器更容易估计学生状态；若只有稀疏完成行为，状态更新会更不可靠。第三，\textbf{奖励可解释性}：多数公开数据集记录的是对错或完成行为，而不是干预后的真实知识增益，因此 RL-LPR 往往需要用 KT/KES 预测的掌握度变化近似奖励。第四，\textbf{结构与语义信息}：知识关系、题目文本和难度信息决定了能否研究图增强、难度感知和 LLM/RAG 增强的路径推荐。

\subsection{习题级 KT 与 ER 数据集}

ASSISTments 系列是 KT、认知诊断和练习推荐中使用最广的数据集之一，来自 ASSISTments 在线辅导平台。其 2009、2012 和 2017 版本分别覆盖不同年份或数据挖掘竞赛子集，常被用于检验模型是否能根据历史答题序列预测后续表现\cite{ASSISTmentsData2009,Cheng2026UniER}。对于 LPR，ASSISTments 的价值在于它提供了相对密集的题目级反馈和知识点标注，便于构造“当前状态--推荐习题--答题反馈”的闭环；局限在于，它主要反映数学练习场景，且离线日志并不直接告诉研究者某条推荐路径会造成怎样的真实因果学习收益。

KDD Cup 2010 中的 Algebra2005 和 Bridge2006 同样是教育数据挖掘中的经典基准。它们来自中学代数智能教学系统，包含多步代数题作答日志和知识概念标注，适合研究长序列答题、时序依赖和知识点掌握预测\cite{Cheng2026UniER}。在 LPR 语境下，这类数据能支持基于知识概念的路径规划和学生模拟器训练，但其领域集中在代数，动作空间和知识结构也受到课程系统设计的限制。

Junyi 数据集来自 Junyi Academy，包含大规模学生解题日志、习题元数据以及练习之间的关系信息。CSEAL 使用 Junyi Academy 这类在线教育系统作为重要应用背景，说明其数据天然适合研究“学习者知识水平”和“学习项目知识结构”的耦合\cite{Liu2019CSEAL}。UniER 的统计显示，在其统一基准口径下，Junyi 包含 25,925,992 条交互、247,606 名学生、722 道练习和 41 个知识概念\cite{Cheng2026UniER}。这种“学生多、题目相对少、交互极多”的结构适合检验模型在大规模用户行为上的稳定性，但也要求研究者注意热门题、重复练习和平台路径引导带来的分布偏置。

\subsection{大规模、多行为与挑战赛数据集}

EdNet 是由 Santa AI 教育平台收集的大规模层级教育数据集。原始 EdNet 论文将其描述为包含数十万学习者和超过一亿级交互的层级数据集，记录时间跨度超过两年，行为类型不仅包括答题，还包括内容浏览等学习活动\cite{Choi2020EdNet}。UniER 在其处理口径下报告 EdNet 含 95,293,926 条交互、784,309 名学生、13,169 道练习和 188 个知识概念\cite{Cheng2026UniER}。因此，EdNet 更适合研究大规模、长历史、多行为的学生建模，但如果任务目标是精确评估某条推荐路径的教学效果，仍然需要额外的学生模拟器或在线实验支撑。

NIPS34 源自 NeurIPS 2020 Education Challenge 所使用的 Eedi 诊断题数据。该挑战关注学生在数学诊断题上的作答预测、问题质量和个性化序列等任务，强调诊断题、技能标签和泛化能力\cite{Wang2020NeurIPSEducation}。UniER 中的 NIPS34 子集包含 1,382,727 条交互、4,918 名学生、948 道练习和 57 个知识概念\cite{Cheng2026UniER}。这类数据适合检验模型在题目稀疏、诊断性强的场景中的泛化能力，也适合讨论 LPR 是否能为学生安排更有信息量的诊断路径。

XES3G5M 是较新的 K-12 在线数学学习数据集，UniER 将其纳入统一 benchmark，并指出其除标准答题序列外，还包含层级知识概念和题目文本等辅助信息\cite{Cheng2026UniER}。这类数据对 LLM/RAG 增强 LPR 尤其重要，因为如果没有题目文本、知识点层级或语义元数据，大模型只能围绕 ID 和稀疏标签进行解释，难以发挥语义理解优势。

\begin{table*}[t]
\centering
\caption{LPR/ER/KT 中常见公开数据集及其建模意义。统计量采用 UniER 附录的统一处理口径。}
\label{tab:lpr-datasets}
\small
\begin{tabularx}{\textwidth}{p{2.55cm}p{2.15cm}p{2.75cm}p{3.0cm}X}
\toprule
数据集 & 场景来源 & 规模（交互/学生/练习/KC） & 主要可用信息 & 对 LPR 的意义与限制 \\
\midrule
ASSISTments 2009 & 在线数学辅导 & 346,860 / 4,217 / 26,688 / 123 & 答题结果、知识点、题目级反馈 & 适合 KT 与短路径练习推荐；但学习收益多需由模拟器近似\cite{ASSISTmentsData2009,Cheng2026UniER}。\\
ASSISTments 2012 & 在线数学辅导 & 6,123,270 / 46,674 / 179,999 / 265 & 大规模题目级答题日志 & 规模更大，适合检验稀疏题目空间下的推荐；跨版本可比性需谨慎\cite{Cheng2026UniER}。\\
ASSISTments 2017 & 数据挖掘竞赛子集 & 942,816 / 686 / 3,162 / 102 & 竞赛式答题序列与 KC 标签 & 学生数较少但序列密集，常用于鲁棒性和路径级 ER 对比\cite{Cheng2026UniER}。\\
Algebra2005 & KDD Cup 2010 代数教学系统 & 813,661 / 575 / 1,084 / 112 & 多步代数题、KC 标注、时序日志 & 适合长序列与代数领域知识建模；领域和课程结构较窄\cite{Cheng2026UniER}。\\
Bridge2006 & KDD Cup 2010 Bridge to Algebra & 3,686,871 / 1,146 / 19,258 / 493 & 大规模代数交互和 KC 标注 & 有利于检验复杂时序动态；路径推荐结果可能受特定课程系统约束\cite{Cheng2026UniER}。\\
NIPS34 & Eedi/NeurIPS 2020 教育挑战 & 1,382,727 / 4,918 / 948 / 57 & 诊断题、技能标签、学生作答 & 适合检验诊断性路径和题目稀疏场景泛化\cite{Wang2020NeurIPSEducation,Cheng2026UniER}。\\
Junyi & Junyi Academy & 25,925,992 / 247,606 / 722 / 41 & 解题日志、习题元数据、练习关系 & 适合大规模学生状态建模和知识结构利用；题目数相对较少，需注意重复练习偏置\cite{Liu2019CSEAL,Cheng2026UniER}。\\
EdNet & Santa AI 在线辅导平台 & 95,293,926 / 784,309 / 13,169 / 188 & 答题、浏览等多行为层级日志 & 适合长历史、多行为建模；真实路径收益仍需额外评估\cite{Choi2020EdNet,Cheng2026UniER}。\\
XES3G5M & K-12 在线数学平台 & 5,549,635 / 18,066 / 7,652 / 865 & 层级知识概念、题目文本和辅助元数据 & 更适合语义增强和 LLM/RAG-LPR；但新数据集的跨论文可比性仍在形成中\cite{Cheng2026UniER}。\\
\bottomrule
\end{tabularx}
\end{table*}

\subsection{面向 LPR 的数据集选择原则}

数据集选择不应简单追求规模最大，而应与研究问题匹配。若目标是研究 KT 与 LPR 的闭环关系，应优先选择具有密集答题反馈和可靠 KC 标注的数据集，如 ASSISTments、Junyi 或 KDD Cup 系列；若目标是研究长历史、多行为和大规模在线系统，应考虑 EdNet；若目标是研究诊断性路径或冷启动泛化，可考虑 NIPS34；若目标是研究 LLM/RAG 或题目语义增强，则应选择包含题目文本、层级知识概念或元数据的数据集，如 XES3G5M。

更严格地说，数据集应当满足“目标--观测”一致性。若论文声称优化真实学习收益 $\Delta \mathbf{s}$，但数据只观测答题正确性 $y_t$，则需要明确说明
\begin{equation}
\widehat{\Delta \mathbf{s}}_t=f_{\mathrm{KT}}(\mathcal{H}_{t+1})-f_{\mathrm{KT}}(\mathcal{H}_{t})
\end{equation}
只是由学生模型推断出的代理收益，而不是直接观测到的因果学习增益。若论文声称路径满足先修关系，则数据中应存在人工知识图、可验证的 KC 层级，或至少有可从日志中统计检验的依赖关系。若论文声称模型具有语义理解能力，则数据中需要题目文本、解析步骤或知识点描述，否则“语义增强”很可能只是额外参数化的 ID 表征。

因此，公开数据集对 LPR 的贡献主要是提供可复现实验环境，而不是直接证明方法在真实教学中的因果有效性。未来更理想的数据应记录推荐动作本身、学生接受或跳过推荐的行为、路径执行过程、阶段性测验结果和长期保持效果。只有当数据从被动日志逐渐转向主动干预日志时，LPR 才能更可靠地区分“模型预测准确”与“推荐确实促进学习”。

\section{方法对比与评估问题}

\begin{table*}[t]
\centering
\caption{代表性学习路径推荐方法对比。}
\label{tab:method-comparison}
\small
\begin{tabularx}{\textwidth}{p{2.2cm}p{1.6cm}p{2.3cm}p{3.4cm}X}
\toprule
\textbf{方法} & \textbf{时间} & \textbf{主要范式} & \textbf{核心贡献} & \textbf{主要局限} \\
\midrule
本体/知识图/多约束方法 & 2007--2018 & 规则与组合优化 & 利用语义、本体、先修关系和约束生成可解释路径 & 状态动态性弱，通常不直接优化长期学习收益 \\
CSEAL & 2019 & MDP + Actor-Critic & 结合认知结构、状态追踪和策略优化 & 状态估计、动作约束和奖励设计仍较依赖特定设置 \\
GEHRL & 2023 & 图增强层级 RL & 分解知识点与学习项目决策，降低动作空间 & 层级训练复杂，低层动作质量影响整体收益 \\
SRC & 2023 & Set-to-Sequence + PG & 从目标集合生成完整概念路径，引入辅助 KT & 实时反馈重规划能力有限 \\
DLPR & 2024 & 难度感知分层规划 & 将题目难度显式纳入路径推荐 & 难度以外的遗忘、情感和多行为因素仍不足 \\
GEPKSD & 2024 & 状态蒸馏 + RL & 用特权知识状态缓解部分可观测问题 & 依赖训练阶段额外状态信息 \\
LIGHT & 2025 & 知识拓扑序列优化 & 融合显式先修和隐式协同关系 & 更多关注图和序列结构，在线交互验证仍有限 \\
KnowLP & 2025 & GraphRAG + 双图 + RL & 自动构建双知识结构图，降低人工建图成本 & 自动图质量验证和噪声控制仍是关键 \\
UNO & 2026 & 离线训练范式 & 面向离线日志和稀疏奖励改进训练 & 离线评估与真实在线效果仍需进一步连接 \\
IB-GRPO & 2026 & LLM + 多指标对齐 & 将 LLM 路径生成与教育指标对齐 & 目前仍以预印本和模拟评估为主，真实教学验证不足 \\
\bottomrule
\end{tabularx}
\end{table*}

\subsection{数据集集中与指标不统一}

上一节说明，ASSISTments、Junyi Academy、KDD Cup、NIPS34、EdNet、XES3G5M 等数据集为 LPR 和 KT 提供了可复现基准。近期 UniER 尝试统一 item-level 和 path-level exercise recommendation 基准，也说明该领域正在意识到数据集割裂问题\cite{Cheng2026UniER}。但是，数据集数量增加并不等于评估问题已经解决，目前仍存在三个困难：
\begin{enumerate}
  \item \textbf{离线指标替代真实学习效果。} AUC、ACC、NDCG 等指标方便计算，但未必等价于真实学习收益。
  \item \textbf{路径质量指标不统一。} 路径合理性、完成率、难度平滑性、知识覆盖率、重复率和长期保持率在不同论文中定义不一。
  \item \textbf{在线实验稀缺。} CSEAL 等少数工作包含在线验证，多数后续方法仍主要依赖离线模拟。
\end{enumerate}

更严格的评价应至少包含三层指标：
\begin{equation}
\text{Prediction} \rightarrow \text{Path Quality} \rightarrow \text{Learning Outcome}.
\end{equation}
第一层评价 KT 或学生模拟器是否可信；第二层评价路径是否满足先修、难度和多样性约束；第三层评价真实学习收益、保持率和学习者体验。

\subsection{过程型路径与细粒度习题推荐}

除主流 RL 方法外，过程挖掘和 DKT 结合的 process-type LPR 也值得关注\cite{Zhang2024ProcessDKT}。它将学习路径视为可包含分支的过程模型，而不是单一路径序列。这与真实学习更接近：学习者可能因反馈不同进入不同分支。未来 LPR 可能从固定序列推荐进一步转向策略图、条件路径或教学流程推荐。

\section{开放问题与未来方向}

\subsection{知识结构自动构建与验证}

知识结构是 LPR 的骨架。传统人工知识图可靠但成本高，GraphRAG 和 LLM 自动建图可扩展但容易引入噪声。一个更稳健的方向是构建“生成--验证--修正”闭环：先由 LLM/RAG 生成候选边，再利用学习日志中的条件依赖验证。例如，如果 $k_a$ 是 $k_b$ 的先修，则未掌握 $k_a$ 的学习者对 $k_b$ 的表现通常受限。可用统计关系近似检验：
\begin{equation}
P(r_{b}=1 \mid s_a < \tau) \ll P(r_{b}=1 \mid s_a \ge \tau).
\end{equation}
该检验不能直接证明因果关系，但可作为自动建图的过滤信号。

\subsection{奖励函数与多目标优化}

强化学习式 LPR 的核心不是简单套用 DQN、PPO 或 Actor-Critic，而是如何设计教育上合理的奖励。单一终局奖励稀疏，单步正确率奖励可能诱导短视策略，难度奖励又可能忽视长期迁移。更合理的奖励应同时包含目标达成、连续掌握提升、难度适配、重复控制和长期保持：
\begin{equation}
R_t=
\alpha\Delta U_t+
\beta D_t-
\lambda C_t-
\rho O_t+
\zeta H_t,
\end{equation}
其中 $\Delta U_t$ 是掌握提升，$D_t$ 是难度适配，$C_t$ 是路径成本，$O_t$ 是冗余或违反先修的惩罚，$H_t$ 是保持或巩固收益。关键问题是这些权重不能只靠经验调参，还需要与教育目标、数据规模和学习阶段匹配。

\subsection{KT--LPR 联合训练}

许多工作把 KT 作为预训练状态估计器，再训练路径策略。这种模块化方式稳定，但误差会层层传播。端到端联合训练可能更优，但也更难解释和调试。未来可考虑三任务联合框架：KT 负责状态追踪，认知诊断负责能力解释，LPR 负责动作决策。三者共享知识表示但保留各自目标函数：
\begin{equation}
\mathcal{L}=
\mathcal{L}_{KT}
+\lambda_{CD}\mathcal{L}_{CD}
+\lambda_{RL}\mathcal{L}_{RL}.
\end{equation}
这种设计有望兼顾预测精度、诊断解释和路径收益。

\subsection{在线验证、公平性与可解释性}

真实教学环境中的 LPR 还必须回答三个问题。第一，离线模拟得到的收益能否迁移到真实在线学习？第二，系统是否对低基础学习者、高基础学习者或不同学习风格群体产生不公平推荐？第三，教师和学生能否理解推荐理由？尤其是深度 RL 和 LLM 方案，若不能给出可信解释，实际部署会受到限制。

可解释 LPR 不应只输出“推荐该题，因为模型分数高”，而应说明推荐与学生状态、目标、先修关系和难度之间的关系。例如：
\begin{quote}
推荐知识点 $k_b$，因为目标集合中 $k_b$ 尚未达标，且其先修 $k_a$ 已达到掌握阈值；推荐习题 $i$，因为其难度接近当前能力估计，并覆盖 $k_b$ 的关键子技能。
\end{quote}
这类解释需要从模型内部状态和外部知识结构中生成，并接受事实一致性检查。

\section{结论}

学习路径推荐已经从早期的规则、知识图和组合优化，逐步发展为融合知识追踪、图神经网络、强化学习和大语言模型的复杂序列决策问题。CSEAL 的关键意义在于把认知结构与策略学习结合起来，使 LPR 从静态路径规划走向动态自适应决策。此后，层级 RL、set-to-sequence、难度感知、知识拓扑优化、状态蒸馏、离线训练和 GraphRAG 自动建图不断扩展该范式。与此同时，KT 的进展为 LPR 提供了更精确的状态估计，但 KT 指标与路径学习收益之间仍存在间隙。一个值得继续推进的思路是把知识点规划、习题执行和状态更新闭合到逐题循环中，用混合过程奖励缓解稀疏反馈，并用局部执行器提升当前知识点下的选题质量。

未来 LPR 研究应避免只追求模型复杂度，而应围绕教育目标建立更可靠的闭环：高质量知识结构、可信学生状态、合理奖励函数、统一评估基准、在线验证和可解释反馈。只有当推荐路径能够在真实学习中稳定提升长期学习效果，并让教师和学生理解其依据时，LPR 才能从算法研究走向可部署的智能教育系统。

\printbibliography[title={参考文献}]

\end{document}
